动态规划先定义状态所包含的信息，再写转移与边界。删除次数、库存、已选区间数等限制必须进入状态；不可达状态应与合法的零收益状态区分。计数题还要保证各条转移对应互不重叠的方案，否则会重复计数。

选择算法时，应证明可利用的结构：二分依赖单调的排除方向，贪心依赖交换或支配关系，分治需要说明跨子问题的贡献如何计算。复杂度按所有状态、转移和数据结构操作累计；空间压缩时还须检查更新顺序是否误用本轮结果。随机算法应另行区分分布正确性、期望成本与最坏成本。
\par\smallskip
